%%%% ijcai26.tex

\typeout{IJCAI--ECAI 26 Instructions for Authors}

% These are the instructions for authors for IJCAI--ECAI 26.

\documentclass{article}
\pdfpagewidth=8.5in
\pdfpageheight=11in

% The file ijcai26.sty is a copy from ijcai22.sty
% The file ijcai22.sty is NOT the same as previous years'
\usepackage{ijcai26}

% Use the postscript times font!
\usepackage{times}
\usepackage{soul}
\usepackage{url}
\usepackage[hidelinks]{hyperref}
\usepackage[utf8]{inputenc}
\usepackage[small]{caption}
\usepackage{graphicx}
\usepackage{amsmath}
\usepackage{amsthm}
\usepackage{booktabs}
\usepackage{algorithm}
\usepackage{algorithmic}
\usepackage[switch]{lineno}

% Comment out this line in the camera-ready submission
\linenumbers

\urlstyle{same}

% the following package is optional:
%\usepackage{latexsym}

% See https://www.overleaf.com/learn/latex/theorems_and_proofs
% for a nice explanation of how to define new theorems, but keep
% in mind that the amsthm package is already included in this
% template and that you must *not* alter the styling.
\newtheorem{example}{Example}
\newtheorem{theorem}{Theorem}
\newtheorem{definition}{Definition}

\newtheorem{subdefinitioninner}{Definition}[definition]
\renewcommand{\thesubdefinitioninner}{\thedefinition.\arabic{subdefinitioninner}}

\newenvironment{subdefinition}
  {\begin{list}{}{\leftmargin=2em}
   \item\begin{subdefinitioninner}}
  {\end{subdefinitioninner}\end{list}}


% Following comment is from ijcai97-submit.tex:
% The preparation of these files was supported by Schlumberger Palo Alto
% Research, AT\&T Bell Laboratories, and Morgan Kaufmann Publishers.
% Shirley Jowell, of Morgan Kaufmann Publishers, and Peter F.
% Patel-Schneider, of AT\&T Bell Laboratories collaborated on their
% preparation.

% These instructions can be modified and used in other conferences as long
% as credit to the authors and supporting agencies is retained, this notice
% is not changed, and further modification or reuse is not restricted.
% Neither Shirley Jowell nor Peter F. Patel-Schneider can be listed as
% contacts for providing assistance without their prior permission.

% To use for other conferences, change references to files and the
% conference appropriate and use other authors, contacts, publishers, and
% organizations.
% Also change the deadline and address for returning papers and the length and
% page charge instructions.
% Put where the files are available in the appropriate places.


% PDF Info Is REQUIRED.

% Please leave this \pdfinfo block untouched both for the submission and
% Camera Ready Copy. Do not include Title and Author information in the pdfinfo section
\pdfinfo{
/TemplateVersion (IJCAI.2026.0)
}

\title{Submodular Maximization with Optimal Hints}

% Single author syntax
\author{
    Author Name
    \affiliations
    Affiliation
    \emails
    email@example.com
}

% Multiple author syntax (remove the single-author syntax above and the \iffalse ... \fi here)
\iffalse
\author{
First Author$^1$
\and
Second Author$^2$\and
Third Author$^{2,3}$\And
Fourth Author$^4$\\
\affiliations
$^1$First Affiliation\\
$^2$Second Affiliation\\
$^3$Third Affiliation\\
$^4$Fourth Affiliation\\
\emails
\{first, second\}@example.com,
third@other.example.com,
fourth@example.com
}
\fi

\begin{document}

\maketitle

\begin{abstract}
    Submodular maximization under cardinality constraints is a fundamental problem with broad applications in several areas such as machine learning, data summarization, feature selection, and influence maximization. The classical framework assumes no prior knowledge about the possible solution sets, and the greedy algorithm provides the best possible approximation $(1-1/e)$ in this setting. We consider the submodular maximization problem with the prior knowledge of a partial optimal solution; we refer to this knowledge as optimal hints.  

There are two possible scenarios with cardinality constraint $\leq k$: 
one where some $\ell$ elements of an optimal solution are given and another where  the "best" $\ell$ elements of some optimal solution is given. We generalize it later as the \emph{goodness} level of the prediction. These scenarios are motivated by the idea that prior knowledge may become available due to domain knowledge from experts,  due to some more expensive prior computations, or as predictions of a learning algorithm. 

We show that the greedy strategy admits to 
$(1-1/e)f(OPT) + (1/e) f(OPT_{\ell})$-approximation for the first case and to $1 - ((k-\ell)/ek) (1 - 1/(k-\ell))^{-\ell}f(OPT)$-approximation for the second case. In addition to theoretical guarantees, we present preliminary experiments characterizing the situations where application of the greedy strategy, in practice, can lead to a close to optimal solution  in the presence of prior knowledge of partial optimal solutions. 
\end{abstract}

\section{Introduction}

We consider the \emph{maximization of a nonnegative submodular function} problem. That is, given submodular function $f: 2^{V} \rightarrow \mathbf{R}_+$ find a set $S \subseteq V$ that maximizes $f(S)$. 

A function $f: 2^{V} \rightarrow \mathbf{R}$ is submodular if for all $A, B \subseteq V$, $f(A) + f(B) \geq  f(A \cup B) + f(A \cap B)$. 

\subsection{Background}
Submodular maximization appears in many areas of computer science, especially combinatorial optimization, max cut, set covering, matroid optimization, sensor placing, and others [references]. A good example is the \emph{max cover under cardinality constraint} problem; where we are asked to choose a collection of sets of size $k$ whose union is maximized. The greedy approach to this problem is well-known to provide a $(1-1/e)$-approximation, and this is optimal under the $P \neq NP$ assumption. In the latter example, and many other settings, the submodular functions of interest are said to be \emph{monotone}, i.e. for any $A\subseteq B \subseteq V, f(A) \leq f(B)$. Interestingly, the \emph{max cover under cardinality constraint} problem can be represented as a maximization of a monotone submodular function under a cardinality constraint; that is, $max_{S \subseteq V} \{f(S) : |S| \leq K\}$. For the cases where this is not trivial, the submodular maximization problem under cardinality constraint is proved to be NP-hard [reference]. 

Here we will consider the case where the functions are \emph{monotone}. For this instance, Nemhauser et al. [Reference] proposed a greedy algorithm with a near-optimal approximation factor of $(1 -\frac{1}{e})$; later on [Reference], it was proven that such factor is the best we can do with only a polynomial number of function queries. However, the greedy algorithm does not perform as well under other constraints like the knapsack constraint.

We will also assume to have access to the submodular function by querying an \emph{oracle}. That is, we have access to the value of $f(S)$ for any $S \subseteq V$ through an \emph{oracle}.

\subsection{Preliminaries}

\section{Contributions}

For any instance $f, V, k$ of the Monotone Submodular Maximization problem and a prediction set $S_l$ where $|S_l| = l \leq k$, and where $T_l = S_l \cup \{g_{l+1}, g_{l+2}, ..., g_k\}$ is the solution of choosing the remaining $k - l$ elements greedily. We present two theorems for the results of $T_l$:

\begin{theorem}[Given no \emph{Goodness} guarantee for the hints]
    $$f(T_l) \geq \Big(1 - \frac{1}{e}\Big)f(OPT) + \frac{1}{e}f(S_l)$$
\end{theorem}

\begin{proof}
    Let $g: 2 ^ V \rightarrow R_+$ and $f: 2^V \rightarrow R_+$ be two submodular functions. Let $f(S) = g(S | S_l)$, and $G_l = \{g_{l+1}, g_{l+2}, ..., g_{k}\}$. Then, from [Reference] we know that:

    $$
    \begin{aligned}
        f(G_l) \geq & \Big(1 - \frac{1}{e}\Big) f(OPT) \\
        g(G_l | S_l) \geq & \Big(1 - \frac{1}{e}\Big) g(OPT | S_l) \\
        g(G_l \cup S_l) - g(S_l) \geq & \Big(1 - \frac{1}{e}\Big) (g(OPT \cup S_l) - g(S_l)) \\
        g(T_l) - g(S_l) \geq & \Big(1 - \frac{1}{e}\Big) g(OPT \cup S_l) - g(S_l) + \frac{1}{e}g(S_l) \\
        g(T_l) \geq & \Big(1 - \frac{1}{e}\Big) g(OPT \cup S_l)  + \frac{1}{e}g(S_l) 
    \end{aligned}      
    $$
\end{proof}

\begin{theorem}[Given some \emph{Goodness} guarantee for the hints]
     If $f(S_l) \geq (1 - (1 - \frac{1}{k})^l + \alpha) f(OPT)$, then
     $$f(T_l) \geq  \Big(1 - \frac{1}{e} + \Big(1 - \frac{1}{k}\Big)^{k-l}\alpha\Big)f(OPT)$$
\end{theorem}

\begin{proof}
    
    Let  $\lambda_i = OPT - f(T_i)$, and $\lambda_0 = f(OPT) - f(S_l)$ then we have:
    
    $$\begin{aligned}
    f(OPT)  &\leq f(OPT \cup T_i)\ \\%&& \text{Monotonicity}\\
            &\leq f(T_i) + \sum_{j = 1}^{k}f(o_j|T_i) \\
            &\leq f(T_i) + \sum_{j = 1}^k f(g_{i+1 - l}|T_i) \\%&& \text{Submodularity and }  g_{l+i+1} = \text{best for the i+1-th iteration} \\
            &= f(T_i) + k f(g_{l+i+1} | T_i) \\
            &= f(T_i) + k (f(T_{i + 1}) - f(T_i)) \\
        f(OPT) - f(T_i) &\leq  k (f(T_{i + 1}) - f(T_i)) \\
        \lambda_i &\leq  k (\lambda_{i} - \lambda_{i + 1}) \\
        \lambda_{i + 1} &\leq  \Big(1 - \frac{1}{k}\Big) \lambda_{i} \\
        \lambda_{k - l} &\leq \Big(1 - \frac{1}{k}\Big) ^ {k - l} \lambda_{0} \\
        &\leq \Big(1 - \frac{1}{k}\Big) ^ {k - l} (f(OPT) - f(S_l)) \\
        &\leq \Big(1 - \frac{1}{k}\Big) ^ {k - l} \Big[f(OPT)  \\
        & - \Big(1 - \Big(1 - \frac{1}{k}\Big)^l + \alpha\Big)f(OPT)\Big] \\
        &\leq \Big(1 - \frac{1}{k}\Big) ^ {k - l}\Big(\Big(1 - \frac{1}{k}\Big)^{l} - \alpha\Big)f(OPT) \\
        &\leq \Big(\Big(1 - \frac{1}{k}\Big)^{k} - \Big(1 - \frac{1}{k}\Big)^{k-l}\alpha\Big)f(OPT) \\
    f(OPT) - f(T_{k-l})    &\leq \Big(\Big(1 - \frac{1}{k}\Big)^{k} - \Big(1 - \frac{1}{k}\Big)^{k-l}\alpha\Big)f(OPT) \\
    \end{aligned}$$

Thus,
    $$\begin{aligned}
    f(T_{k-l})    &\geq \Big(1 -\Big(1 - \frac{1}{k}\Big)^{k} + \Big(1 - \frac{1}{k}\Big)^{k-l}\alpha\Big)f(OPT) \\
                &\geq \Big(1 -\frac{1}{e} + \Big(1 - \frac{1}{k}\Big)^{k-l}\alpha\Big)f(OPT)\\
    \end{aligned}$$

\end{proof}

We also present a framework for this study that can be used to quantify the general case of having a predictor that gives hints at any moment of the well known greedy algorithm, and an example of such type of algorithm.

\subsection{Framework}

\begin{definition}
    \begin{subdefinition}
        Let $S_l = \{g_1, g_2, ..., g_l\}$ where $g_t$ is the element chosen in the $t$-th step of our algorithm. 
    \end{subdefinition}
    \begin{subdefinition}
        Let $O_l = \{o^{'}_{l}, o^{'}_{l + 1}, ..., o^{'}_{k}\}$: the best $k-l$ elements to build on $S_l$
    \end{subdefinition}
    \begin{subdefinition}
        $\mathrm{OPT}_l = S_{l-1} \cup O_l$: the optimal solution given an initial partial solution $S_l$.
    \end{subdefinition}
    \begin{subdefinition}
            $\Gamma_l = OPT_l - f(S_{l-1})$
    \end{subdefinition}
\end{definition}




    

\begin{definition}[Predictor-quality-on-correctness]
    For every step $i$ there exists a constant $\rho_i>0$ such that

    $$\mathbf{E}\big[\Delta(p_i\mid S_{i-1}) \mid p_i\in O_i\big]\;\ge\; \rho_i\cdot\frac{\mathrm{OPT_i}-f(S_{i-1})}{k}.$$
\end{definition}




\textbf{Theorem 3.} Given an SMP, $S_l$ (as above), a predictor that will give a prediction $p_i$ at step $l < i \leq k$ and $Pr[p_i \in O_i] = 1$, then ...:


\textbf{Proof:}

Similar to Yanhui's analysis the proof goes along the same lines on step A through C, changing $OPT$ with $OPT_i$ and $O$ with $O_i$. Moreover we have,

\[
c_i = (1 - \beta_i) +\beta_i\alpha_i\rho_i \tag{*}
\]

\[
\mathbf{E}\big[f(S_i) - f(S_{i-1})\big] \geq \frac{c_i}{k} \Gamma_i \tag{**}
\]

Then, by definition we have $\Gamma_{i+1} = (\mathrm{OPT}_{i+1} - \mathrm{OPT}_{i}) +  \Gamma_i - (f(S_i) - f(S_{i-1}))$. Using (**) we have: 

\[
\begin{aligned}
\mathbf{E}\big[\Gamma_{i+1}\big] &= (\mathrm{OPT}_{i+1} - \mathrm{OPT}_{i}) + \Gamma_i - \mathbf{E}\big[f(S_i) - f(S_{i-1})\big] \\
    &\leq (\mathrm{OPT}_{i+1} - \mathrm{OPT}_{i}) + \Gamma_i - \frac{c_i}{k}\Gamma_i & (\mathrm{OPT_{i+1} - OPT_i} \leq 0) \\
    &\leq \Big(1 - \frac{c_i}{k}\Big)\Gamma_i
\end{aligned}
\]

% The {\it IJCAI--ECAI 26 Proceedings} will be printed from electronic
% manuscripts submitted by the authors. These must be PDF ({\em Portable
%         Document Format}) files formatted for 8-1/2$''$ $\times$ 11$''$ paper.

% \subsection{Length of Papers}


% All paper {\em submissions} to the main track must have a maximum of seven pages, plus at most two for references / acknowledgements / contribution statement / ethics statement.

% The length rules may change for final camera-ready versions of accepted papers and
% differ between tracks. Some tracks may disallow any contents other than the references in the last two pages, whereas others allow for any content in all pages. Similarly, some tracks allow you to buy a few extra pages should you want to, whereas others don't.

% If your paper is accepted, please carefully read the notifications you receive, and check the proceedings submission information website\footnote{\url{https://proceedings.ijcai.org/info}} to know how many pages you can use for your final version. That website holds the most up-to-date information regarding paper length limits at all times.


% \subsection{Word Processing Software}

% As detailed below, IJCAI has prepared and made available a set of
% \LaTeX{} macros and a Microsoft Word template for use in formatting
% your paper. If you are using some other word processing software, please follow the format instructions given below and ensure that your final paper looks as much like this sample as possible.

% \section{Style and Format}

% \LaTeX{} and Word style files that implement these instructions
% can be retrieved electronically. (See Section~\ref{stylefiles} for
% instructions on how to obtain these files.)

% \subsection{Layout}

% Print manuscripts two columns to a page, in the manner in which these
% instructions are printed. The exact dimensions for pages are:
% \begin{itemize}
%     \item left and right margins: .75$''$
%     \item column width: 3.375$''$
%     \item gap between columns: .25$''$
%     \item top margin---first page: 1.375$''$
%     \item top margin---other pages: .75$''$
%     \item bottom margin: 1.25$''$
%     \item column height---first page: 6.625$''$
%     \item column height---other pages: 9$''$
% \end{itemize}

% All measurements assume an 8-1/2$''$ $\times$ 11$''$ page size. For
% A4-size paper, use the given top and left margins, column width,
% height, and gap, and modify the bottom and right margins as necessary.

% \subsection{Format of Electronic Manuscript}

% For the production of the electronic manuscript, you must use Adobe's
% {\em Portable Document Format} (PDF). A PDF file can be generated, for
% instance, on Unix systems using {\tt ps2pdf} or on Windows systems
% using Adobe's Distiller. There is also a website with free software
% and conversion services: \url{http://www.ps2pdf.com}. For reasons of
% uniformity, use of Adobe's {\em Times Roman} font is strongly suggested.
% In \LaTeX2e{} this is accomplished by writing
% \begin{quote}
%     \mbox{\tt $\backslash$usepackage\{times\}}
% \end{quote}
% in the preamble.\footnote{You may want to also use the package {\tt
%             latexsym}, which defines all symbols known from the old \LaTeX{}
%     version.}

% Additionally, it is of utmost importance to specify the {\bf
%         letter} format (corresponding to 8-1/2$''$ $\times$ 11$''$) when
% formatting the paper. When working with {\tt dvips}, for instance, one
% should specify {\tt -t letter}.

% \subsection{Papers Submitted for Review vs. Camera-ready Papers}
% In this document, we distinguish between papers submitted for review (henceforth, submissions) and camera-ready versions, i.e., accepted papers that will be included in the conference proceedings. The present document provides information to be used by both types of papers (submissions / camera-ready). There are relevant differences between the two versions. Find them next.

% \subsubsection{Anonymity}
% For the main track and some of the special tracks, submissions must be anonymous; for other special tracks they must be non-anonymous. The camera-ready versions for all tracks are non-anonymous. When preparing your submission, please check the track-specific instructions regarding anonymity.

% \subsubsection{Submissions}
% The following instructions apply to submissions:
% \begin{itemize}
% \item If your track requires submissions to be anonymous, they must be fully anonymized as discussed in the Modifications for Blind Review subsection below; in this case, Acknowledgements and Contribution Statement sections are not allowed.

% \item If your track requires non-anonymous submissions, you should provide all author information at the time of submission, just as for camera-ready papers (see below); Acknowledgements and Contribution Statement sections are allowed, but optional.

% \item Submissions must include line numbers to facilitate feedback in the review process . Enable line numbers by uncommenting the command {\tt \textbackslash{}linenumbers} in the preamble.

% \item The limit on the number of  content pages is \emph{strict}. All papers exceeding the limits will be desk rejected.
% \end{itemize}

% \subsubsection{Camera-Ready Papers}
% The following instructions apply to camera-ready papers:

% \begin{itemize}
% \item Authors and affiliations are mandatory. Explicit self-references are allowed. It is strictly forbidden to add authors not declared at submission time.

% \item Acknowledgements and Contribution Statement sections are allowed, but optional.

% \item Line numbering must be disabled. To achieve this, comment or disable {\tt \textbackslash{}linenumbers} in the preamble.

% \item For some of the tracks, you can exceed the page limit by purchasing extra pages.
% \end{itemize}

% \subsection{Title and Author Information}

% Center the title on the entire width of the page in a 14-point bold
% font. The title must be capitalized using Title Case. For non-anonymous papers, author names and affiliations should appear below the title. Center author name(s) in 12-point bold font. On the following line(s) place the affiliations.

% \subsubsection{Author Names}

% Each author name must be followed by:
% \begin{itemize}
%     \item A newline {\tt \textbackslash{}\textbackslash{}} command for the last author.
%     \item An {\tt \textbackslash{}And} command for the second to last author.
%     \item An {\tt \textbackslash{}and} command for the other authors.
% \end{itemize}

% \subsubsection{Affiliations}

% After all authors, start the affiliations section by using the {\tt \textbackslash{}affiliations} command.
% Each affiliation must be terminated by a newline {\tt \textbackslash{}\textbackslash{}} command. Make sure that you include the newline after the last affiliation, too.

% \subsubsection{Mapping Authors to Affiliations}

% If some scenarios, the affiliation of each author is clear without any further indication (\emph{e.g.}, all authors share the same affiliation, all authors have a single and different affiliation). In these situations you don't need to do anything special.

% In more complex scenarios you will have to clearly indicate the affiliation(s) for each author. This is done by using numeric math superscripts {\tt \$\{\^{}$i,j, \ldots$\}\$}. You must use numbers, not symbols, because those are reserved for footnotes in this section (should you need them). Check the authors definition in this example for reference.

% \subsubsection{Emails}

% This section is optional, and can be omitted entirely if you prefer. If you want to include e-mails, you should either include all authors' e-mails or just the contact author(s)' ones.

% Start the e-mails section with the {\tt \textbackslash{}emails} command. After that, write all emails you want to include separated by a comma and a space, following the order used for the authors (\emph{i.e.}, the first e-mail should correspond to the first author, the second e-mail to the second author and so on).

% You may ``contract" consecutive e-mails on the same domain as shown in this example (write the users' part within curly brackets, followed by the domain name). Only e-mails of the exact same domain may be contracted. For instance, you cannot contract ``person@example.com" and ``other@test.example.com" because the domains are different.


% \subsubsection{Modifications for Blind Review}
% When submitting to a track that requires anonymous submissions,
% in order to make blind reviewing possible, authors must omit their
% names, affiliations and e-mails. In place
% of names, affiliations and e-mails, you can optionally provide the submission number and/or
% a list of content areas. When referring to one's own work,
% use the third person rather than the
% first person. For example, say, ``Previously,
% Gottlob~\shortcite{gottlob:nonmon} has shown that\ldots'', rather
% than, ``In our previous work~\cite{gottlob:nonmon}, we have shown
% that\ldots'' Try to avoid including any information in the body of the
% paper or references that would identify the authors or their
% institutions, such as acknowledgements. Such information can be added post-acceptance to be included in the camera-ready
% version.
% Please also make sure that your paper metadata does not reveal
% the authors' identities.

% \subsection{Abstract}

% Place the abstract at the beginning of the first column 3$''$ from the
% top of the page, unless that does not leave enough room for the title
% and author information. Use a slightly smaller width than in the body
% of the paper. Head the abstract with ``Abstract'' centered above the
% body of the abstract in a 12-point bold font. The body of the abstract
% should be in the same font as the body of the paper.

% The abstract should be a concise, one-paragraph summary describing the
% general thesis and conclusion of your paper. A reader should be able
% to learn the purpose of the paper and the reason for its importance
% from the abstract. The abstract should be no more than 200 words long.

% \subsection{Text}

% The main body of the text immediately follows the abstract. Use
% 10-point type in a clear, readable font with 1-point leading (10 on
% 11).

% Indent when starting a new paragraph, except after major headings.

% \subsection{Headings and Sections}

% When necessary, headings should be used to separate major sections of
% your paper. (These instructions use many headings to demonstrate their
% appearance; your paper should have fewer headings.). All headings should be capitalized using Title Case.

% \subsubsection{Section Headings}

% Print section headings in 12-point bold type in the style shown in
% these instructions. Leave a blank space of approximately 10 points
% above and 4 points below section headings.  Number sections with
% Arabic numerals.

% \subsubsection{Subsection Headings}

% Print subsection headings in 11-point bold type. Leave a blank space
% of approximately 8 points above and 3 points below subsection
% headings. Number subsections with the section number and the
% subsection number (in Arabic numerals) separated by a
% period.

% \subsubsection{Subsubsection Headings}

% Print subsubsection headings in 10-point bold type. Leave a blank
% space of approximately 6 points above subsubsection headings. Do not
% number subsubsections.

% \paragraph{Titled paragraphs.} You should use titled paragraphs if and
% only if the title covers exactly one paragraph. Such paragraphs should be
% separated from the preceding content by at least 3pt, and no more than
% 6pt. The title should be in 10pt bold font and to end with a period.
% After that, a 1em horizontal space should follow the title before
% the paragraph's text.

% In \LaTeX{} titled paragraphs should be typeset using
% \begin{quote}
%     {\tt \textbackslash{}paragraph\{Title.\} text} .
% \end{quote}

% \subsection{Special Sections}

% \subsubsection{Appendices}
% You may move some of the contents of the paper into one or more appendices that appear after the main content, but before references. These appendices count towards the page limit and are distinct from the supplementary material that can be submitted separately through CMT. Such appendices are useful if you would like to include highly technical material (such as a lengthy calculation) that will disrupt the flow of the paper. They can be included both in papers submitted for review and in camera-ready versions; in the latter case, they will be included in the proceedings (whereas the supplementary materials will not be included in the proceedings).
% Appendices are optional. Appendices must appear after the main content.
% Appendix sections must use letters instead of Arabic numerals. In \LaTeX, you can use the {\tt \textbackslash{}appendix} command to achieve this followed by  {\tt \textbackslash section\{Appendix\}} for your appendix sections.

% \subsubsection{Ethical Statement}

% Ethical Statement is optional. You may include an Ethical Statement to discuss  the ethical aspects and implications of your research. The section should be titled \emph{Ethical Statement} and be typeset like any regular section but without being numbered. This section may be placed on the References pages.

% Use
% \begin{quote}
%     {\tt \textbackslash{}section*\{Ethical Statement\}}
% \end{quote}

% \subsubsection{Acknowledgements}

% Acknowledgements are optional. In the camera-ready version you may include an unnumbered acknowledgments section, including acknowledgments of help from colleagues, financial support, and permission to publish. This is not allowed in the anonymous submission. If present, acknowledgements must be in a dedicated, unnumbered section appearing after all regular sections but before references.  This section may be placed on the References pages.

% Use
% \begin{quote}
%     {\tt \textbackslash{}section*\{Acknowledgements\}}
% \end{quote}
% to typeset the acknowledgements section in \LaTeX{}.


% \subsubsection{Contribution Statement}

% Contribution Statement is optional. In the camera-ready version you may include an unnumbered Contribution Statement section, explicitly describing the contribution of each of the co-authors to the paper. This is not allowed in the anonymous submission. If present, Contribution Statement must be in a dedicated, unnumbered section appearing after all regular sections but before references.  This section may be placed on the References pages.

% Use
% \begin{quote}
%     {\tt \textbackslash{}section*\{Contribution Statement\}}
% \end{quote}
% to typeset the Contribution Statement section in \LaTeX{}.

% \subsubsection{References}

% The references section is headed ``References'', printed in the same
% style as a section heading but without a number. A sample list of
% references is given at the end of these instructions. Use a consistent
% format for references. The reference list should not include publicly unavailable work.

% \subsubsection{Order of Sections}
% Sections should be arranged in the following order:
% \begin{enumerate}
%     \item Main content sections (numbered)
%     \item Appendices (optional, numbered using capital letters)
%     \item Ethical statement (optional, unnumbered)
%     \item Acknowledgements (optional, unnumbered)
%     \item Contribution statement (optional, unnumbered)
%     \item References (required, unnumbered)
% \end{enumerate}

% \subsection{Citations}

% Citations within the text should include the author's last name and
% the year of publication, for example~\cite{gottlob:nonmon}.  Append
% lowercase letters to the year in cases of ambiguity.  Treat multiple
% authors as in the following examples:~\cite{abelson-et-al:scheme}
% or~\cite{bgf:Lixto} (for more than two authors) and
% \cite{brachman-schmolze:kl-one} (for two authors).  If the author
% portion of a citation is obvious, omit it, e.g.,
% Nebel~\shortcite{nebel:jair-2000}.  Collapse multiple citations as
% follows:~\cite{gls:hypertrees,levesque:functional-foundations}.
% \nocite{abelson-et-al:scheme}
% \nocite{bgf:Lixto}
% \nocite{brachman-schmolze:kl-one}
% \nocite{gottlob:nonmon}
% \nocite{gls:hypertrees}
% \nocite{levesque:functional-foundations}
% \nocite{levesque:belief}
% \nocite{nebel:jair-2000}

% \subsection{Footnotes}

% Place footnotes at the bottom of the page in a 9-point font.  Refer to
% them with superscript numbers.\footnote{This is how your footnotes
%     should appear.} Separate them from the text by a short
% line.\footnote{Note the line separating these footnotes from the
%     text.} Avoid footnotes as much as possible; they interrupt the flow of
% the text.

% \section{Illustrations}

% Place all illustrations (figures, drawings, tables, and photographs)
% throughout the paper at the places where they are first discussed,
% rather than at the end of the paper.

% They should be floated to the top (preferred) or bottom of the page,
% unless they are an integral part
% of your narrative flow. When placed at the bottom or top of
% a page, illustrations may run across both columns, but not when they
% appear inline.

% Illustrations must be rendered electronically or scanned and placed
% directly in your document. They should be cropped outside \LaTeX{},
% otherwise portions of the image could reappear during the post-processing of your paper.
% When possible, generate your illustrations in a vector format.
% When using bitmaps, please use 300dpi resolution at least.
% All illustrations should be understandable when printed in black and
% white, albeit you can use colors to enhance them. Line weights should
% be 1/2-point or thicker. Avoid screens and superimposing type on
% patterns, as these effects may not reproduce well.

% Number illustrations sequentially. Use references of the following
% form: Figure 1, Table 2, etc. Place illustration numbers and captions
% under illustrations. Leave a margin of 1/4-inch around the area
% covered by the illustration and caption.  Use 9-point type for
% captions, labels, and other text in illustrations. Captions should always appear below the illustration.

% \section{Tables}

% Tables are treated as illustrations containing data. Therefore, they should also appear floated to the top (preferably) or bottom of the page, and with the captions below them.

% \begin{table}
%     \centering
%     \begin{tabular}{lll}
%         \hline
%         Scenario  & $\delta$ & Runtime \\
%         \hline
%         Paris     & 0.1s     & 13.65ms \\
%         Paris     & 0.2s     & 0.01ms  \\
%         New York  & 0.1s     & 92.50ms \\
%         Singapore & 0.1s     & 33.33ms \\
%         Singapore & 0.2s     & 23.01ms \\
%         \hline
%     \end{tabular}
%     \caption{Latex default table}
%     \label{tab:plain}
% \end{table}

% \begin{table}
%     \centering
%     \begin{tabular}{lrr}
%         \toprule
%         Scenario  & $\delta$ (s) & Runtime (ms) \\
%         \midrule
%         Paris     & 0.1          & 13.65        \\
%                   & 0.2          & 0.01         \\
%         New York  & 0.1          & 92.50        \\
%         Singapore & 0.1          & 33.33        \\
%                   & 0.2          & 23.01        \\
%         \bottomrule
%     \end{tabular}
%     \caption{Booktabs table}
%     \label{tab:booktabs}
% \end{table}

% If you are using \LaTeX, you should use the {\tt booktabs} package, because it produces tables that are better than the standard ones. Compare Tables~\ref{tab:plain} and~\ref{tab:booktabs}. The latter is clearly more readable for three reasons:

% \begin{enumerate}
%     \item The styling is better thanks to using the {\tt booktabs} rulers instead of the default ones.
%     \item Numeric columns are right-aligned, making it easier to compare the numbers. Make sure to also right-align the corresponding headers, and to use the same precision for all numbers.
%     \item We avoid unnecessary repetition, both between lines (no need to repeat the scenario name in this case) as well as in the content (units can be shown in the column header).
% \end{enumerate}

% \section{Formulas}

% IJCAI's two-column format makes it difficult to typeset long formulas. A usual temptation is to reduce the size of the formula by using the {\tt small} or {\tt tiny} sizes. This doesn't work correctly with the current \LaTeX{} versions, breaking the line spacing of the preceding paragraphs and title, as well as the equation number sizes. The following equation demonstrates the effects (notice that this entire paragraph looks badly formatted, and the line numbers no longer match the text):
% %
% \begin{tiny}
%     \begin{equation}
%         x = \prod_{i=1}^n \sum_{j=1}^n j_i + \prod_{i=1}^n \sum_{j=1}^n i_j + \prod_{i=1}^n \sum_{j=1}^n j_i + \prod_{i=1}^n \sum_{j=1}^n i_j + \prod_{i=1}^n \sum_{j=1}^n j_i
%     \end{equation}
% \end{tiny}%

% Reducing formula sizes this way is strictly forbidden. We {\bf strongly} recommend authors to split formulas in multiple lines when they don't fit in a single line. This is the easiest approach to typeset those formulas and provides the most readable output%
% %
% \begin{align}
%     x = & \prod_{i=1}^n \sum_{j=1}^n j_i + \prod_{i=1}^n \sum_{j=1}^n i_j + \prod_{i=1}^n \sum_{j=1}^n j_i + \prod_{i=1}^n \sum_{j=1}^n i_j + \nonumber \\
%     +   & \prod_{i=1}^n \sum_{j=1}^n j_i.
% \end{align}%

% If a line is just slightly longer than the column width, you may use the {\tt resizebox} environment on that equation. The result looks better and doesn't interfere with the paragraph's line spacing: %
% \begin{equation}
%     \resizebox{.91\linewidth}{!}{$
%             \displaystyle
%             x = \prod_{i=1}^n \sum_{j=1}^n j_i + \prod_{i=1}^n \sum_{j=1}^n i_j + \prod_{i=1}^n \sum_{j=1}^n j_i + \prod_{i=1}^n \sum_{j=1}^n i_j + \prod_{i=1}^n \sum_{j=1}^n j_i
%         $}.
% \end{equation}%

% This last solution may have to be adapted if you use different equation environments, but it can generally be made to work. Please notice that in any case:

% \begin{itemize}
%     \item Equation numbers must be in the same font and size as the main text (10pt).
%     \item Your formula's main symbols should not be smaller than {\small small} text (9pt).
% \end{itemize}

% For instance, the formula
% %
% \begin{equation}
%     \resizebox{.91\linewidth}{!}{$
%             \displaystyle
%             x = \prod_{i=1}^n \sum_{j=1}^n j_i + \prod_{i=1}^n \sum_{j=1}^n i_j + \prod_{i=1}^n \sum_{j=1}^n j_i + \prod_{i=1}^n \sum_{j=1}^n i_j + \prod_{i=1}^n \sum_{j=1}^n j_i + \prod_{i=1}^n \sum_{j=1}^n i_j
%         $}
% \end{equation}
% %
% would not be acceptable because the text is too small.

% \section{Examples, Definitions, Theorems and Similar}

% Examples, definitions, theorems, corollaries and similar must be written in their own paragraph. The paragraph must be separated by at least 2pt and no more than 5pt from the preceding and succeeding paragraphs. They must begin with the kind of item written in 10pt bold font followed by their number (e.g.: {\bf Theorem 1}),
% optionally followed by a title/summary between parentheses in non-bold font and ended with a period (in bold).
% After that the main body of the item follows, written in 10 pt italics font (see below for examples).

% In \LaTeX{} we strongly recommend that you define environments for your examples, definitions, propositions, lemmas, corollaries and similar. This can be done in your \LaTeX{} preamble using \texttt{\textbackslash{newtheorem}} -- see the source of this document for examples. Numbering for these items must be global, not per-section (e.g.: Theorem 1 instead of Theorem 6.1).

% \begin{example}[How to write an example]
%     Examples should be written using the example environment defined in this template.
% \end{example}

% \begin{theorem}
%     This is an example of an untitled theorem.
% \end{theorem}

% You may also include a title or description using these environments as shown in the following theorem.

% \begin{theorem}[A titled theorem]
%     This is an example of a titled theorem.
% \end{theorem}

% \section{Proofs}

% Proofs must be written in their own paragraph(s) separated by at least 2pt and no more than 5pt from the preceding and succeeding paragraphs. Proof paragraphs should start with the keyword ``Proof." in 10pt italics font. After that the proof follows in regular 10pt font. At the end of the proof, an unfilled square symbol (qed) marks the end of the proof.

% In \LaTeX{} proofs should be typeset using the \texttt{\textbackslash{proof}} environment.

% \begin{proof}
%     This paragraph is an example of how a proof looks like using the \texttt{\textbackslash{proof}} environment.
% \end{proof}


% \section{Algorithms and Listings}

% Algorithms and listings are a special kind of figures. Like all illustrations, they should appear floated to the top (preferably) or bottom of the page. However, their caption should appear in the header, left-justified and enclosed between horizontal lines, as shown in Algorithm~\ref{alg:algorithm}. The algorithm body should be terminated with another horizontal line. It is up to the authors to decide whether to show line numbers or not, how to format comments, etc.

% In \LaTeX{} algorithms may be typeset using the {\tt algorithm} and {\tt algorithmic} packages, but you can also use one of the many other packages for the task.

% \begin{algorithm}[tb]
%     \caption{Example algorithm}
%     \label{alg:algorithm}
%     \textbf{Input}: Your algorithm's input\\
%     \textbf{Parameter}: Optional list of parameters\\
%     \textbf{Output}: Your algorithm's output
%     \begin{algorithmic}[1] %[1] enables line numbers
%         \STATE Let $t=0$.
%         \WHILE{condition}
%         \STATE Do some action.
%         \IF {conditional}
%         \STATE Perform task A.
%         \ELSE
%         \STATE Perform task B.
%         \ENDIF
%         \ENDWHILE
%         \STATE \textbf{return} solution
%     \end{algorithmic}
% \end{algorithm}

% \section{\LaTeX{} and Word Style Files}\label{stylefiles}

% The \LaTeX{} and Word style files are available on the IJCAI--ECAI 26
% website, \url{https://2026.ijcai.org/}.
% These style files implement the formatting instructions in this
% document.

% The \LaTeX{} files are {\tt ijcai26.sty} and {\tt ijcai26.tex}, and
% the Bib\TeX{} files are {\tt named.bst} and {\tt ijcai26.bib}. The
% \LaTeX{} style file is for version 2e of \LaTeX{}, and the Bib\TeX{}
% style file is for version 0.99c of Bib\TeX{} ({\em not} version
% 0.98i). 

% The Microsoft Word style file consists of a single file, {\tt ijcai26.docx}. 
% %This template differs from the one used for IJCAI--23.

% These Microsoft Word and \LaTeX{} files contain the source of the
% present document and may serve as a formatting sample.

% Further information on using these styles for the preparation of
% papers for IJCAI--ECAI 26 can be obtained by contacting {\tt
%         proceedings@ijcai.org}.

% \appendix

% \section*{Ethical Statement}

% There are no ethical issues.

% \section*{Acknowledgments}

% The preparation of these instructions and the \LaTeX{} and Bib\TeX{}
% files that implement them was supported by Schlumberger Palo Alto
% Research, AT\&T Bell Laboratories, and Morgan Kaufmann Publishers.
% Preparation of the Microsoft Word file was supported by IJCAI.  An
% early version of this document was created by Shirley Jowell and Peter
% F. Patel-Schneider.  It was subsequently modified by Jennifer
% Ballentine, Thomas Dean, Bernhard Nebel, Daniel Pagenstecher,
% Kurt Steinkraus, Toby Walsh, Carles Sierra, Marc Pujol-Gonzalez,
% Francisco Cruz-Mencia and Edith Elkind.


%% The file named.bst is a bibliography style file for BibTeX 0.99c
\bibliographystyle{named}
\bibliography{ijcai26}

\end{document}

